\chapter{リスク管理}
\label{ch:risk}

ヘッジファンドが潰れる理由は、予測が外れたことよりも、リスクの取り方を誤ったことのほうが多い。この章では、リスクを測る道具（リスクモデル、VaR）と、リスクを管理する仕組みをまとめる。

\section{リスクの種類}

\begin{table}[H]
\centering\small
\caption{ヘッジファンドが直面するリスク}
\label{tab:risks}
\begin{tabular}{@{}P{0.2\linewidth}P{0.72\linewidth}@{}}
\toprule
リスク & 中身 \\
\midrule
市場リスク & 価格・金利・為替・変動率が不利に動く。 \\
ファクターのリスク & 割安・モメンタムなど、持っているスタイルがまとめて損をする。 \\
流動性のリスク & 売りたいときに売れない、売ると価格が大きく動く（市場の流動性）。 \\
資金繰りのリスク & 借入れの担保の追加を求められる、貸し手が資金を引き上げる、投資家が解約する（資金の流動性）。 \\
信用・取引相手のリスク & 債券の発行体やデリバティブの相手、プライムブローカーが破綻する。 \\
モデルのリスク & モデルの前提（相関、分布、流動性）が崩れる。 \\
混雑のリスク & 同じ取引をしている人が一斉に手じまう。 \\
オペレーションのリスク & システムの誤り、発注の誤り、不正。 \\
\bottomrule
\end{tabular}
\end{table}

市場の流動性と資金の流動性は、互いに強め合う（Brunnermeier・Pedersen 2009 \cite{brunnermeierpedersen2009}）。価格が下がると担保の価値が下がり、担保の追加を求められて売る。売ると価格がさらに下がる。この「流動性のらせん」が、ヘッジファンドの大きな損失の多くの背景にある（第 II 部、第 \ref{ch:failures} 章）。

\section{リスクの物差し}

\subsection{ボラティリティ}

最も基本的な物差しは、リターンの標準偏差（ボラティリティ）である。ポートフォリオのボラティリティは $\sigma_p = \sqrt{\bm w\T\Sigma\bm w}$。

\subsection{VaR}

\textbf{VaR}（バリュー・アット・リスク）は、「ある期間に、ある確率（信頼水準）で超えない損失」である。例えば「1 日の 99\% の VaR が 1 億円」とは、100 日に 1 日程度、1 億円を超える損失が出るという意味である。

\begin{itemize}
\item \textbf{正規分布を仮定}：$\mathrm{VaR}_\alpha = z_\alpha\,\sigma_p\,\sqrt{h} \times A$（$A$ は資産額、$z_{0.99} \approx 2.33$、$h$ は日数）。
\item \textbf{過去のデータ（ヒストリカル）}：過去の各日のリターンを、いまのポートフォリオに当てはめ、損失の分布の分位点をとる。
\item \textbf{モンテカルロ}：リスクの要因を確率モデルで多数生成し、損失の分布を作る。オプションのような非線形な資産に向く。
\end{itemize}

\subsection{期待ショートフォール}

VaR は「どれだけ悪いことが起きうるか」の境目しか見ず、その先の損失の大きさを無視する。\textbf{期待ショートフォール}（ES、CVaR）は、VaR を超えた日の損失の平均である。正規分布なら
\begin{equation}
\mathrm{ES}_\alpha = \sigma_p\,\frac{\phi(z_\alpha)}{1-\alpha}
\end{equation}
（$\phi$ は標準正規の密度）。裾の厚い分布では、ES のほうがリスクをよく表す。銀行の規制も、VaR から ES に移った。

\begin{exbox}[VaR の計算]
資産 100 億円、年率のボラティリティ 10\% のポートフォリオ。1 日のボラティリティは $10\% / \sqrt{252} \approx 0.63\%$。
\begin{itemize}
\item 1 日の 99\% の VaR：$2.33 \times 0.63\% \times 100 \approx 1.47$ 億円。
\item 1 日の 99\% の ES：$0.63\% \times \phi(2.33)/0.01 \approx 0.63\% \times 2.67 \approx 1.68$ 億円。
\item 10 日の 99\% の VaR：$1.47 \times \sqrt{10} \approx 4.6$ 億円（独立を仮定）。
\end{itemize}
実際の株のリターンは裾が厚く、99\% を超える損失は、正規分布の想定よりずっと大きくなりやすい。
\end{exbox}

\subsection{ストレス・テスト}

過去の危機（1987 年のブラックマンデー、1998 年、2008 年、2020 年 3 月）や、仮想の状況（金利が 2\% 上がり、株が 20\% 下がる）を、いまのポートフォリオに当てはめて損失を計算する。統計の物差しでは見えない「相関が崩れたとき」の損失を見るのに欠かせない。

\section{リスクモデル}

\subsection{なぜ必要か}

$N$ 銘柄の共分散 $\Sigma$ には、$N(N+1)/2$ 個の数がある。3,000 銘柄なら約 450 万個。日次のデータが 1 年（252 日）しかなければ、標本の共分散は特異（逆行列がない）になり、推定の誤差が非常に大きい。そこで、構造を仮定して推定する数を減らす。

\subsection{ファクター型のリスクモデル}

銘柄のリターンを、少数（数十）の共通のファクターと、銘柄固有の部分に分ける。
\begin{equation}
\bm r_t = B\,\bm f_t + \bm u_t, \qquad \Sigma = B F B\T + D
\end{equation}
$B$（$N \times K$）は各銘柄のファクターへの感応度、$F$（$K \times K$）はファクターの共分散、$D$ は銘柄固有の分散の対角行列。推定する数は $NK + K(K+1)/2 + N$ に減る。

\begin{itemize}
\item \textbf{ファンダメンタル型}（Barra 型）：感応度 $B$ を、業種のダミー、規模・割安・モメンタム・ボラティリティなどの銘柄の特性で直接与え、毎期の横断面の回帰でファクターのリターン $\bm f_t$ を推定する。商用のリスクモデルの多くはこの形。
\item \textbf{統計型}：主成分分析でファクターを取り出す。解釈は難しいが、データだけで作れる。
\item \textbf{マクロ型}：金利・インフレ・景気などの経済変数をファクターにし、時系列の回帰で感応度を推定する。
\end{itemize}

ファクター型のモデルは、ポートフォリオのリスクを「どのファクターから来ているか」に分解できる。マルチマネージャーが、運用チームにファクターのリスクを打ち消させるときに使うのも、これである。

\subsection{縮小推定}

Ledoit・Wolf（2004）\cite{ledoitwolf2004} は、標本の共分散 $S$ を、構造を持った目標 $F$（すべての相関が等しい行列など）に向けて縮めた
\begin{equation}
\hat\Sigma = \alpha F + (1-\alpha) S
\end{equation}
を提案し、平均二乗誤差を最小にする $\alpha$ をデータから計算する式を示した。計算が簡単で頑健なので、広く使われている。LLM-BLM の BLM でも、直近 252 日の Ledoit--Wolf の推定を使った。

\subsection{ランダム行列による雑音の除去}

$N$ 銘柄、$T$ 日の互いに独立な雑音から作った相関行列の固有値は、$q = N/T$ のとき、Marchenko--Pastur 分布に従い、ほぼ
\begin{equation}
\lambda_\pm = \big(1 \pm \sqrt{q}\big)^2
\end{equation}
の範囲に収まる。実際の株の相関行列の固有値のうち、この範囲を超える少数（市場全体、業種）だけが本当の構造で、範囲の中は雑音と見なせる。範囲の中の固有値を平均値に置き換えることで、相関行列を「掃除」する（Laloux ほか 1999 \cite{laloux1999}、Plerou ほか 1999 \cite{plerou1999} の観察に基づく）。

\section{ボラティリティの予測}

リターンの平均の予測は難しいが、ボラティリティは予測できる。大きな変動の後には大きな変動が続く（ボラティリティのクラスタリング）からである。

\subsection{EWMA}

指数加重の移動平均で分散を更新する。
\begin{equation}
\hat\sigma_t^2 = \lambda\,\hat\sigma_{t-1}^2 + (1-\lambda)\,r_{t-1}^2
\end{equation}
RiskMetrics（1994）は日次で $\lambda = 0.94$ を使った。

\subsection{GARCH}

Engle（1982）\cite{engle1982} の ARCH を一般化した Bollerslev（1986）\cite{bollerslev1986} の GARCH(1,1) は
\begin{equation}
\sigma_t^2 = \omega + a\,r_{t-1}^2 + b\,\sigma_{t-1}^2, \qquad a + b < 1
\end{equation}
で、長期の平均 $\omega/(1-a-b)$ に戻る力を持つ。株では、下げた後に変動が大きくなる非対称（レバレッジ効果）を入れた GJR-GARCH や EGARCH もよく使う。

\subsection{実現ボラティリティ}

日中の 5 分ごとなどの高頻度のリターンの二乗和で、その日の分散を直接測る（実現分散）。Corsi（2009）\cite{corsi2009} の HAR モデルは、日・週・月の実現分散の平均の線形回帰で、翌日の実現分散を予測する。簡単だが強い。

\section{リスクの管理の仕組み}

\subsection{リスクの上限}

運用会社は、次のような上限を置き、超えたら機械的にポジションを減らす。

\begin{itemize}
\item ポートフォリオ全体の VaR・ボラティリティの上限。
\item グロス・ネットのレバレッジの上限。
\item 1 銘柄・1 業種・1 国への集中の上限。
\item ファクターへの感応度の上限。
\item 売買高に対するポジションの大きさの上限（例：1 日の売買高の数日分で手じまえる大きさ）。
\item 損失の上限（ストップロス、ドローダウンの上限）。
\end{itemize}

\subsection{ドローダウンの管理}

運用チームの損失が一定に達したら資金の割り当てを減らし、さらに損をしたらポジションを閉じる、という規則は、マルチマネージャーの標準である。これは、損失が続く運用者を早く止める（腕がないか、市場が変わった可能性が高い）とともに、会社全体の損失を限る。一方で、多くのファンドが同じ規則で同時に損切りすると、損切りの売りが市場を動かす。

\subsection{リスクの予算}

全体のリスク（例：年率 10\% のボラティリティ）を、戦略・運用チーム・ファクターに配分する。期待される情報レシオの高いところに多く配る。リスクパリティ（\chapref{ch:pc}）は、期待を使わずに等しく配る極端な場合である。

\section{リスクモデルの限界}

\begin{itemize}
\item \textbf{危機には相関が上がる}：普段は低い相関の資産が、危機には一斉に同じ方向に動く。分散の効果が、最も必要なときに消える。
\item \textbf{過去の窓の長さ}：過去 1 年の平穏な期間で推定したリスクは、次の危機のリスクを過小評価する。危機の直後は逆に過大評価する。
\item \textbf{混雑は見えない}：ほかの運用者が何を持っているかは、価格のデータからは分かりにくい。
\item \textbf{流動性は見えない}：平常時の売買高は、危機のときに売れる量を表さない。
\end{itemize}

\begin{projbox}
LLM-BLM のポートフォリオは、ロングのみで、事前の TE（リスクモデルで予測した、比べる相手との差のボラティリティ）に上限を置いた。小型株の S1 の追記で、事前の TE の上限を年 3\% のまま「校正しない」としたのは、小型株では事前の TE と実現の TE の比を測った過去がないからである。リスクモデルの予測した TE と、実際に実現した TE を比べることは、それ自体がリスクモデルの検証になる。
\end{projbox}
